% Electrical-design fragment for insertion after the Task subsection.

\subsection{Electrical Design and Low-Level Control}
\label{sec:electrical_design}

The robot has two actuation domains. The differential-drive base uses brushed DC gear motors and a dedicated drive controller. The 3-DoF manipulator uses Hiwonder serial-bus servos controlled through an ESP32. Both are supplied from one 12~V lithium-polymer battery using a star power distribution. The drive controller is connected directly to the battery through its barrel jack, so the intermittent motor current follows a short, low-resistance path. A 5~V/3~A buck converter provides a regulated supply for the Raspberry Pi~4 and the USB-powered devices, including the eye-in-hand camera and the ESP32 controller. This separates the logic and sensing supply from motor voltage drops and switching noise. Heavy-gauge wiring and a current-rated connector are used on the motor branch, while lower-current wiring is used for the sensing and signal branch.

\subsubsection{Mobile-base drive electronics and sensing}
\label{sec:electrical_drive}

Two TorqueNADO permanent-magnet DC (PMDC) gear motors with 60:1 gearboxes and output-shaft quadrature encoders drive the base. They are driven through a controller that includes a microcontroller, a dual full-bridge stage, and per-channel sensing. The drive stage is rated at about 10~A continuous per channel over a 12--18~V supply range, with current sensing up to about 30~A peak. It receives PWM velocity commands, reads the encoders, and provides voltage and current telemetry. Through this controller the base exposes one external interface: a velocity command in, with position and electrical feedback out. The encoders provide 360 counts per revolution, or 1440 counts per revolution with four-fold quadrature decoding. This gives approximately $0.125^\circ$ per count at the output shaft. The controller reads the electrical state using an internal voltage divider and a Hall-effect current sensor for each channel. Both measurements are sampled by a 10-bit ADC and feed the motor model and observer described in Section~\ref{sec:electrical_observer}.

\subsubsection{DC-motor model, state-space model, and velocity observer}
\label{sec:electrical_observer}

Each drive motor is represented as a PMDC motor. The parameter values used in the model are $R=16~\Omega$, $L=1~\mathrm{mH}$, $K_e=K_t=0.013178$ ($K_e$ in $\mathrm{V\,s/rad}$, $K_t$ in $\mathrm{N\,m/A}$), $J=10^{-3}~\mathrm{kg\,m^2}$, and $b/J=0.8819~\mathrm{s}^{-1}$:
\begin{equation}
\label{eq:elec_motor}
V=R\,i+L\dot i+K_e\,\omega,\qquad
J\dot\omega=K_t\,i-b\,\omega,\qquad
\dot\theta=\omega .
\end{equation}

With the state vector $\mathbf{x}=[\theta\ \omega\ i]^\top$ and the input $u=V$, the model is written as $\dot{\mathbf{x}}=\mathbf{A}\mathbf{x}+\mathbf{B}u$ and $\mathbf{y}=\mathbf{C}\mathbf{x}$. Only the encoder position $\theta$ and motor current $i$ are measured:
\begin{equation}
\label{eq:elec_state_space}
\mathbf{A}=
\begin{bmatrix}
0&1&0\\
0&-b/J&K_t/J\\
0&-K_e/L&-R/L
\end{bmatrix}
=
\begin{bmatrix}
0&1&0\\
0&-0.882&13.18\\
0&-13.18&-16000
\end{bmatrix},
\qquad
\mathbf{B}=
\begin{bmatrix}
0\\
0\\
1/L
\end{bmatrix}
=
\begin{bmatrix}
0\\
0\\
1000
\end{bmatrix},
\end{equation}
where $\mathbf{C}=[\,\mathbf{e}_1\ \ \mathbf{e}_3\,]^\top$ selects the measured position and current.

The electrical time constant is $\tau_e=L/R=62.5~\mu\mathrm{s}$, while the mechanical time constant is $\tau=J/b=1.13~\mathrm{s}$. Current therefore settles much faster than the mechanical speed. The wheel controller still needs the unmeasured speed $\omega$. A direct finite difference,
$\omega_{\mathrm{diff}}[k]=(\theta[k]-\theta[k-1])/T_s$, amplifies encoder quantization and measurement noise, especially at low speed. A minimum-order observer is used instead. Since two of the three states are measured, only $\omega$ is estimated.

Let $\mathbf{x}_a=[\theta\ i]^\top$ denote the measured states and $x_b=\omega$ the unmeasured state. Using the auxiliary state $\eta=\hat{\omega}-\mathbf{K}_o\mathbf{x}_a$ gives
\begin{equation}
\label{eq:elec_observer}
\begin{gathered}
\dot{\hat{\eta}}=\hat{A}\,\hat{\eta}
+\hat{\mathbf{B}}\,\mathbf{x}_a+\hat{F}\,u,
\qquad
\hat{\omega}=\hat{\eta}+\mathbf{K}_o\mathbf{x}_a,\\
\hat{A}=A_{bb}-\mathbf{K}_o\mathbf{A}_{ab},
\qquad
\hat{\mathbf{B}}=\hat{A}\mathbf{K}_o+\mathbf{A}_{ba}
-\mathbf{K}_o\mathbf{A}_{aa},
\qquad
\hat{F}=B_b-\mathbf{K}_o\mathbf{B}_a .
\end{gathered}
\end{equation}

The observer error follows $\dot e_b=(A_{bb}-\mathbf{K}_o\mathbf{A}_{ab})e_b$, so its single pole is set directly through $\mathbf{K}_o$. Placing this pole at $\mu=-20$ gives the gain $\mathbf{K}_o=[\,0.1095\ \ {-}1.4425\,]$, and hence $\hat{A}=-20$, $\hat{\mathbf{B}}=[-2.19\ \ {-}2.30\times10^{4}]$, and $\hat{F}=1442.5$. The corresponding observer time constant is $50~\mathrm{ms}$, with an approximate $0.25~\mathrm{s}$ settling time. The estimate avoids the noise amplification caused by direct differentiation.

\subsubsection{Discrete implementation and speed control}
\label{sec:electrical_control}

Estimation and speed control run at $T_s=20~\mathrm{ms}$, or 50~Hz. The observer in \eqref{eq:elec_observer} is discretized using forward Euler:
\[
\hat{\eta}[k+1]=\hat{\eta}[k]+T_s\left(
\hat{A}\hat{\eta}[k]+\hat{\mathbf{B}}\mathbf{x}_a[k]+\hat{F}V[k]\right).
\]
The armature-current measurement is affected by high-frequency PWM and electromagnetic noise, so it is filtered with a first-order exponential moving average of cut-off frequency $f_c$; the wheel speed is then regulated by a discrete PID law acting on the speed error:
\begin{equation}
\label{eq:elec_pid}
i_{\mathrm{filt}}[k]=\alpha i_{\mathrm{raw}}[k]
+(1-\alpha)i_{\mathrm{filt}}[k-1],
\qquad
f_c=\frac{\alpha f_s}{2\pi(1-\alpha)};
\qquad
u[k]=K_pe[k]+K_iT_s\sum_{j=0}^{k}e[j]
+K_d\frac{e[k]-e[k-1]}{T_s},
\end{equation}
where $e[k]=\omega_{\mathrm{ref}}[k]-\hat{\omega}[k]$. The PWM command is clamped, and the integrator is held when saturation occurs to prevent windup.

The body command is converted to wheel-speed references through the inverse differential-drive relation. For track width $L_t$ and wheel radius $r_d$,
\[
\omega_{R,\mathrm{ref}}=
\frac{v+\tfrac{1}{2}L_t\omega_z}{r_d},
\qquad
\omega_{L,\mathrm{ref}}=
\frac{v-\tfrac{1}{2}L_t\omega_z}{r_d}.
\]
At a small computational cost, the observer yields a lower-noise speed estimate with an explicitly selectable convergence rate.

\subsubsection{Manipulator actuation and embedded bridge}
\label{sec:electrical_arm}

The three active joints use Hiwonder serial-bus servos. The fourth wrist joint is passive and is mechanically leveled. Each active servo includes a local controller, a magnetic encoder, and an internal position loop. The servo receives a target position and a movement time, then reports position, input voltage, and temperature. The three servos share a half-duplex TTL bus at 115200~baud using 8N1 framing and IDs 1--3. A position change of 1000 pulses corresponds to $240^\circ$, giving
\[
\frac{\mathrm{d}p}{\mathrm{d}\theta}
=
\frac{1000}{240}\frac{180}{\pi}
=
238.7324
\quad \text{pulses/rad}.
\]

The ESP32 acts as a bridge between the host computer and the servo bus. A USB serial link carries arm commands from the Raspberry Pi~4 to the ESP32. The ESP32 drives the servo bus through a dedicated hardware UART. Because the bus is half-duplex, the ESP32 switches its transceiver between transmit and receive states before each command and before each reply. The host sends the command vector \texttt{[p1,p2,p3,time\_ms]}. The ESP32 clamps the commands, sends timed moves to the servos, and streams read-back positions at 20~Hz. Each bus frame begins with \texttt{0x55 0x55}, followed by the servo ID, length, command, payload, and checksum.

The host converts between joint angles and servo pulses using
\begin{equation}
\label{eq:elec_joint_map}
p_i=s_i\theta_i+o_i,
\qquad
\theta_i=\frac{p_i-o_i}{s_i},
\qquad
\theta_4=2\pi-\frac{\pi}{2}-\theta_{\mathrm{sh}}-\theta_{\mathrm{el}},
\end{equation}
where the pulse-scale vector is
$s=[238.73,\ {-}238.73,\ 238.73]$~pulses/rad and the offset vector is
$o=[625,\ 875,\ 500]$. The sign of each scale term accounts for the joint mounting direction. The pulse-angle mapping is based on the manufacturer relation and must be checked on the physical arm by commanding known angles, reading the reported pulses, and fitting the scale and offset values.

\subsubsection{Simulation-to-hardware signal path and layered safety}
\label{sec:electrical_sim_to_real}

The command path is the same in simulation and on the physical arm up to the servo bus. A target produced by the reinforcement-learning policy or a planner is expressed in joint radians, converted to pulses using \eqref{eq:elec_joint_map}, clamped, and sent to the ESP32. The ESP32 writes the timed move to the bus, and each servo closes its own position loop. Feedback returns in the opposite direction. Read-back pulses are converted to joint angles, the passive-wrist relation is applied, and the joint state is sent to the host at 20~Hz. If feedback is temporarily unavailable, the bridge echoes the previous command, so the joint-state stream remains continuous.

The base path is shorter. The host sends a body-velocity command, the differential-drive relation produces the wheel references, and the base controller regulates the motor speed while returning encoder and electrical telemetry. Safety is implemented in four layers. The host rejects targets outside the allowed range. The ESP32 clamps every pulse to the valid interval. The servos enforce their own angle, voltage, and temperature limits. Finally, the base controller uses current limits to bound motor torque independently of the PWM duty cycle.
